%! AP Statistics — Unit 1, Lesson 1.9 — Group Activity ANSWER KEY
\documentclass[10pt]{article}
\usepackage{apstats-article}
\usepackage{apstats-key}

%=============================================================================
\begin{document}

\pageheader{Unit 1, Lesson 1.9}{Group Activity --- Answer Key}
\begin{tabularx}{\linewidth}{X X X}
Name:\;\ans{Key} & Date:\; & Period:\;
\end{tabularx}

\vspace{0.10in}

\begin{tcolorbox}[colback=sky, colframe=navy, boxrule=0.8pt, arc=2mm,
    left=3mm, right=3mm, top=2mm, bottom=2mm]
\small
\textbf{Part 1} (8 min): Compute five-number summaries and check for outliers.\\
\textbf{Part 2} (7 min): Interpret the pre-drawn parallel boxplots.\\
\textbf{Part 3} (5 min): Write an AP-quality comparison and peer-review.
\end{tcolorbox}

\vspace{0.10in}

{\large\bfseries\color{navy} Part 1 \quad Compute Five-Number Summaries}

\vspace{0.08in}

\noindent\small
\textbf{Athletes} ($n=15$):\enspace
6, 7, 7, 8, 8, 8, 9, \textbf{9}, 9, 10, 10, 10, 11, 11, 12\\[3pt]
\textbf{Non-Athletes} ($n=15$):\enspace
5, 6, 6, 7, 7, 7, 7, \textbf{8}, 8, 8, 9, 9, 10, 10, 14

\vspace{0.10in}

\begin{enumerate}[leftmargin=1.5em, itemsep=6pt]

  \item Median is the \ans{8th} value; $Q_1$ and $Q_3$ are the medians of
        the lower and upper \ans{7} values.

  \item Five-number summary:\\[4pt]
        \renewcommand{\arraystretch}{1.6}
        \begin{tabularx}{\linewidth}{@{} l Y Y Y Y Y @{}}
        \rowcolor{sky}
        \textbf{Group} & $\min$ & $Q_1$ & $M$ & $Q_3$ & $\max$ \\
        \midrule
        Athletes      & \ans{6} & \ans{8} & \ans{9} & \ans{10} & \ans{12} \\
        Non-Athletes  & \ans{5} & \ans{7} & \ans{8} & \ans{9}  & \ans{14} \\
        \end{tabularx}

        \vspace{4pt}
        {\small\textit{Athletes $Q_1$: median of \{6,7,7,8,8,8,9\} $= 8$ (4th of 7).
        $Q_3$: median of \{9,10,10,10,11,11,12\} $= 10$.\\
        Non-Athletes $Q_1$: median of \{5,6,6,7,7,7,7\} $= 7$.
        $Q_3$: median of \{8,8,9,9,10,10,14\} $= 9$.}}

  \item \ans{Athletes IQR $= 10 - 8 = 2$ hr. \quad Non-Athletes IQR $= 9 - 7 = 2$ hr.}\\[4pt]
        Upper fence $= 9 + 1.5(2) = 9 + 3 = $ \ans{12 hr}\\[3pt]
        Is 14 an outlier? \ans{Yes, $14 > 12$ (upper fence).}\\[3pt]
        Last non-outlier $= $ \ans{10 hr}

\end{enumerate}

\vspace{0.10in}

{\large\bfseries\color{navy} Part 2 \quad Interpret the Parallel Boxplots}

\vspace{0.08in}

\begin{center}
\begin{tikzpicture}[scale=1.1]
  \draw[thick] (0.3, 0) -- (6.30, 0);
  \foreach \v/\x in {4/0.00, 5/0.55, 6/1.10, 7/1.65, 8/2.20, 9/2.75,
                      10/3.30, 11/3.85, 12/4.40, 13/4.95, 14/5.50, 15/6.05} {
    \draw (\x, -0.12) -- (\x, 0.12);
    \node[below, font=\small] at (\x, -0.12) {\v};
  }
  \node[below, font=\small\itshape] at (3.02, -0.56) {Hours of sleep per night};
  \node[left, font=\small] at (0.0, 1.10) {Athletes};
  \draw[thick, navy] (1.10, 1.10) -- (2.20, 1.10);
  \draw[thick, navy] (2.20, 0.85) rectangle (3.30, 1.35);
  \draw[thick, navy] (2.75, 0.85) -- (2.75, 1.35);
  \draw[thick, navy] (3.30, 1.10) -- (4.40, 1.10);
  \draw[thick, navy] (1.10, 0.95) -- (1.10, 1.25);
  \draw[thick, navy] (4.40, 0.95) -- (4.40, 1.25);
  \node[left, font=\small] at (0.0, 0.35) {Non-Athletes};
  \draw[thick, navy] (0.55, 0.35) -- (1.65, 0.35);
  \draw[thick, navy] (1.65, 0.10) rectangle (2.75, 0.60);
  \draw[thick, navy] (2.20, 0.10) -- (2.20, 0.60);
  \draw[thick, navy] (2.75, 0.35) -- (3.30, 0.35);
  \draw[thick, navy] (0.55, 0.20) -- (0.55, 0.50);
  \draw[thick, navy] (3.30, 0.20) -- (3.30, 0.50);
  \filldraw[navy] (5.50, 0.35) circle (0.10);
\end{tikzpicture}
\end{center}

\begin{enumerate}[leftmargin=1.5em, itemsep=6pt, resume]

  \item \ans{Yes, the displayed values match the computed five-number summaries.
        The outlier dot at 14 hours represents one non-athlete who slept
        14 hours per night --- a value that exceeds the upper fence of 12 hours.}

  \item Athletes between 8 and 10 hr: $\approx$ \ans{50\%} (the IQR box).\\[3pt]
        Non-Athletes between 7 and 9 hr: $\approx$ \ans{50\%} (the IQR box).

  \item \ans{Athletes have a higher median (9 hr) than Non-Athletes (8 hr),
        a difference of 1 hour.}

  \item \ans{Both groups have equal IQRs of 2 hours ($Q_3 - Q_1$), meaning
        the middle 50\% of each group's sleep times are equally consistent.
        However, Non-Athletes have a higher outlier (14 hr), suggesting at
        least one non-athlete has a very unusual sleep schedule.}

\end{enumerate}

\vspace{0.10in}

{\large\bfseries\color{navy} Part 3 \quad Write an AP-Quality Comparison}

\begin{enumerate}[leftmargin=1.5em, itemsep=6pt, resume]

  \item
  \textbf{Shape:} \ans{Both distributions of nightly sleep hours appear approximately
  symmetric. The athletes' distribution is roughly symmetric with equal whiskers
  (IQR box centered near the median), while the non-athletes' distribution is
  right-skewed due to the high outlier at 14 hours.}\\[5pt]
  \textbf{Outlier(s):} \ans{There is one outlier in the Non-Athletes group at
  14 hours per night. This exceeds the upper fence of $9 + 1.5(2) = 12$ hours
  and is plotted separately as a dot on the modified boxplot. The Athletes group
  has no outliers.}\\[5pt]
  \textbf{Center:} \ans{The median nightly sleep for athletes (9 hours) was
  1 hour \textit{higher than} the median for non-athletes (8 hours), indicating
  that athletes tended to get more sleep per night.}\\[5pt]
  \textbf{Spread:} \ans{The IQR for athletes (2 hours) was \textit{equal to} the
  IQR for non-athletes (2 hours), meaning the middle 50\% of each group's nightly
  sleep times were equally consistent. Excluding the non-athlete outlier, the range
  for athletes ($12 - 6 = 6$ hr) was slightly \textit{greater than} for non-athletes
  ($10 - 5 = 5$ hr).}

  \item Peer review: teacher circulates; credit for naming both groups, comparative word,
        specific values, and both center+spread addressed.

\end{enumerate}

\begin{teachernote}
\textbf{Grading notes:}
\begin{itemize}[itemsep=2pt]
  \item Part 1 (Q1--Q3): 6 pts. Q1: 1 pt (8th, 7). Q2: 2 pts (all 10 values). Q3: 3 pts (IQR both groups, fence, outlier id, last non-outlier).
  \item Part 2 (Q4--Q7): 4 pts (1 each). Q6: any direction word + 1 hr difference.
  \item Part 3 (Q8): SOCS = 4 pts, 1 per component. Each must use comparative language, name both groups, include context. Shape: 1 pt for noting skew in non-athletes due to outlier. Outlier: must give fence calc. Center: comparative, both medians named. Spread: IQR comparison required.
\end{itemize}
\end{teachernote}

\end{document}
